% generated from papers/XIV-equalnorm/paper_XIV_equalnorm.tex -- edit the paper, then rerun assemble.py
\chapter[The equal-norm case]{The equal-norm case: norm-separation is not a hypothesis\\ but a characterization, and the gauge-equivariant repair}\label{ch:XIV}
\chapterorigin{This chapter is Paper XIV of the series (\texttt{papers/XIV-equalnorm/}).}
\begin{summary}
Chapter~\ref{ch:XIII} proved $\cA_{K,S}^\sigma=C(Y_S)$ (its Theorem 1.4) under the
hypothesis that $\Ns:J_S\to\N$ be injective, and asked whether the hypothesis could be
removed. It cannot, and the reason is that it is not a hypothesis at all.

The diagonal is always the fixed-point algebra of the \emph{compact} gauge action:
$C(Y_S)=\cA_{K,S}^{\T^S}$, unconditionally. The modular flow is the restriction of the gauge
action along $\iota(t)_\pp=\Ns\pp^{it}$, so $\cA^\sigma=\cA^{H}$ with
$H=\overline{\iota(\R)}\le\T^S$, and by Kronecker $H=\T^S$ exactly when the $\log\Ns\pp$ are
$\Q$-linearly independent. Since $\Z$-independence of the $\log\Ns\pp$ is literally the
injectivity of $\Ns$ on $J_S$, we obtain
\[
\cA_{K,S}^\sigma=C(Y_S)
\iff \{\log\Ns\pp\}_{\pp\in S}\ \Q\text{-independent}
\iff S\ \text{norm-separated}.
\]
The three conditions are equivalent; there is nothing to remove.

The correct structure in the general case is not the proposed direct sum
$\bigoplus_n\cA^\sigma_n$, which is not a direct sum --- the pieces are neither orthogonal
nor ideals, since $e_\aaa e_\bb=e_{\aaa\vee\bb}\ne0$ and the product of a degree-$n$ and a
degree-$m$ element is again in $\cA^\sigma$. It is a grading. Writing
$L=\{c\in\Z^{(S)}:\sum_\pp c_\pp\log\Ns\pp=0\}$ for the lattice of multiplicative relations
among the norms, $\cA^\sigma$ is $L$-graded with degree-zero part $C(Y_S)$ and
$H=L^\perp$.

Neither of the two proposed intrinsic characterizations of $C(Y_S)$ inside $\cA^\sigma$
works as stated. Property A recovers $C(Y_S)$ as the fixed points of the gauge action, but
the gauge action of $\T^S$ is precisely the datum \emph{not} determined by $(\cA,\sigma)$
when $L\ne0$; the argument is circular. Property B asserts that $C(Y_S)$ is the unique Cartan
subalgebra, which is what was to be proved.

What does work is a change of hypothesis rather than a removal. If $\Phi$ is equivariant for
the full gauge actions --- up to an isomorphism of the acting tori, so that no bijection
$S_1\to S_2$ is presupposed --- then $\Phi(C(Y_{S_1}))=C(Y_{S_2})$ unconditionally, and the
reconstruction of Chapter~\ref{ch:XIII} (its Theorem 2.5) holds for every number field and every prime
set: a norm-preserving bijection $S_1\to S_2$, which turns out to induce the torus
isomorphism after all; the matched valuation fibres, with $G_S$ as a compact space; the
$\mathrm{KMS}_\beta$ simplices, transition loci and $\beta_c$; and, when the symmetry actions
are intertwined, the Frobenius classes modulo inertia and the chain $\{\Xi_\beta\}$. For
$S_i=P_{K_i}$ the norm multisets agree, so $\zeta_{K_1}=\zeta_{K_2}$: the fields are
arithmetically equivalent. The narrow class number is not among the conclusions; it was not
recovered in Chapter~\ref{ch:XIII} either.

Finally we record a limitation of Chapter~\ref{ch:XIII} that we had not noticed there: every $K\ne\Q$
has infinitely many split primes, so $P_K$ is \emph{never} norm-separated. The theorem of
Chapter~\ref{ch:XIII} therefore covers the classical full prime set only for $K=\Q$.
\end{summary}


\section{Norm-separation is a characterization}

We use the notation of the foundational paper [Ch.~\ref{ch:Main}]: $\cA_{K,S}=C(Y_S)\rtimes J_S$ is
spanned by the elements $\mu_\aaa f\mu_\bb^*$, and it is the $C^*$-algebra of the \'etale
groupoid $\cG_S|_{Y_S}$ whose arrows are the pairs $(\aaa\bb^{-1},y)$ with
$y,\,\aaa\bb^{-1}y\in Y_S$, with diagonal $C(Y_S)$ and with the scaling cocycle
$c(\aaa\bb^{-1},y)=\log\Ns\aaa-\log\Ns\bb$ implementing the flow. Chapters~\ref{ch:II}, \ref{ch:III} and \ref{ch:XIII} of
the series, referred to by numeral, are in preparation; every fact used from them is
re-proved in place. The finer, integer-valued cocycle
\[
\widetilde c:\cG_S|_{Y_S}\longrightarrow I_S=\Z^{(S)},\qquad
\widetilde c(\aaa\bb^{-1},y)=\aaa\bb^{-1},
\]
is continuous and locally constant, and composing it with characters defines the
\emph{gauge action} of the compact group $\T^S=\widehat{I_S}$:
\[
\gamma_z(f)(\gamma)=z^{\widetilde c(\gamma)}f(\gamma),\qquad\text{so}\qquad
\gamma_z\bigl(\mu_\pp\bigr)=z_\pp\mu_\pp,\qquad \gamma_z|_{C(Y_S)}=\id,\qquad z\in\T^S,
\]
strongly continuous because $\widetilde c$ is locally constant. Since
$c=\langle\widetilde c,\log\Ns\rangle$, the modular flow is a restriction of the gauge
action: $\sigma_t=\gamma_{\iota(t)}$ with $\iota(t)_\pp=\Ns\pp^{it}$ (the identity through
which Chapter~\ref{ch:III} sees that $\{\sigma_t\}$ has compact closure).

\begin{lemma}[The diagonal is the gauge fixed-point algebra]\label{XIV:lem:gauge}
$\cA_{K,S}^{\T^S}=C(Y_S)$, for every $K$ and every $S$, with no hypothesis.
\end{lemma}

\begin{proof}
$\gamma_z(\mu_\aaa f\mu_\bb^*)=z^{\aaa-\bb}\mu_\aaa f\mu_\bb^*$ in multi-index notation, so a
spectral subspace argument gives $\cA^{\T^S}=\overline{\spn}\{\mu_\aaa f\mu_\aaa^*\}$, and
$\mu_\aaa f\mu_\aaa^*=\alpha_\aaa(f)\in C(Y_S)$.
\end{proof}

\begin{definition}\label{XIV:def:L}
The \emph{relation lattice} of $S$ is
\[
L=L(S):=\Bigl\{c\in\Z^{(S)}:\ \sum_{\pp\in S}c_\pp\log\Ns\pp=0\Bigr\}
=\Bigl\{c:\ \prod_\pp\Ns\pp^{c_\pp}=1\Bigr\},
\]
a subgroup of $\Z^{(S)}$.
\end{definition}

\begin{theorem}[Equivalence]\label{XIV:thm:equiv}
The following are equivalent.
\begin{enumerate}
\item $\Ns:J_S\to\N$ is injective, i.e.\ $S$ is norm-separated;
\item $L(S)=0$, i.e.\ $\{\log\Ns\pp\}_{\pp\in S}$ is $\Q$-linearly independent;
\item $\iota(\R)$ is dense in $\T^S$;
\item $\cA_{K,S}^\sigma=C(Y_S)$.
\end{enumerate}
\end{theorem}

\begin{proof}
$(1)\Leftrightarrow(2)$: $\Ns\aaa=\Ns\bb$ with $\aaa,\bb\in J_S$ means
$\sum(a_\pp-b_\pp)\log\Ns\pp=0$ with $a-b\in\Z^{(S)}$; injectivity says the only such
relation is $a=b$, i.e.\ $L=0$. A $\Q$-linear relation clears denominators to a $\Z$-linear
one, so $\Z$-independence and $\Q$-independence agree.

$(2)\Leftrightarrow(3)$: $\overline{\iota(\R)}$ is a closed subgroup of $\T^S$, and its
annihilator in $\widehat{\T^S}=\Z^{(S)}$ is $\{c:\prod_\pp\Ns\pp^{ic_\pp t}=1\ \forall t\}=L$.
By Pontryagin duality $\overline{\iota(\R)}=\T^S$ iff $L=0$; this is Kronecker's theorem.

$(3)\Leftrightarrow(4)$: $\cA^\sigma=\cA^{\overline{\iota(\R)}}$, since the fixed-point
algebra of a one-parameter group equals that of the closure of its image; combine with
Lemma~\ref{XIV:lem:gauge}. Conversely if $L\ne0$, pick $0\ne c\in L$ and write $c=\aaa-\bb$ with
$\aaa,\bb\in J_S$ coprime; then $\Ns\aaa=\Ns\bb$ and $\mu_\aaa\mu_\bb^*\in\cA^\sigma$ is not
in $C(Y_S)$, the two ideals being distinct.
\end{proof}

\begin{corollary}[Nothing to remove]\label{XIV:cor:nothing}
The hypothesis of Theorem 1.4 of Chapter~\ref{ch:XIII} is not a sufficient condition that might be
weakened: it is equivalent to the conclusion. In the unseparated case $\cA^\sigma$ is strictly
larger than $C(Y_S)$ and non-commutative, and no argument can identify them.
\end{corollary}

\begin{proposition}[$P_K$ is never separated for $K\ne\Q$]\label{XIV:prop:PK}
Let $K\ne\Q$. By the Chebotarev density theorem \cite[Ch.~VII, \S13]{Neukirch}, applied to
the Galois closure of $K$, there are infinitely many rational primes that split completely
in $K$; the $[K:\Q]\ge2$ primes of $K$ above such a $p$ have residue degree $1$ and hence
equal norm $p$. So $P_K$ is not norm-separated, and Theorem 1.4 of Chapter~\ref{ch:XIII} applies to
the full prime set only when $K=\Q$.
\end{proposition}

\begin{remark}[Numerical illustration]\label{XIV:rem:num}
The first split prime is $5$ for $\Q(i)$, $7$ for $\Q(\sqrt{-3})$ and $\Q(\sqrt2)$, and $11$
for $\Q(\sqrt5)$; in each case two primes of norm $p$. This is a limitation of Chapter~\ref{ch:XIII}
that we had not noticed there: the classical case, $S=P_K$ over a general
number field, is precisely the case the theorem does not cover.
\end{remark}

\section{The structure of $\cA^\sigma$ in general}

\begin{proposition}[The proposed stratification is not a direct sum]\label{XIV:prop:notsum}
Set $\cA_n^\sigma=\overline{\spn}\{\mu_\aaa f\mu_\bb^*:\Ns\aaa=\Ns\bb=n\}$. Then
$\cA^\sigma\ne\bigoplus_n\cA_n^\sigma$: the summands are not orthogonal, since
$e_\aaa e_\bb=e_{\mathrm{lcm}(\aaa,\bb)}\ne0$ for $\aaa,\bb$ of different norms, and they are
not ideals, the product of an element of $\cA_n^\sigma$ with one of $\cA_m^\sigma$ lying in
$\cA^\sigma$ but in neither.
\end{proposition}

\begin{proof}
$e_\aaa=\mu_\aaa\mu_\aaa^*\in\cA^\sigma_{\Ns\aaa}$ and $e_\aaa e_\bb=e_{\mathrm{lcm}}\ne0$;
for the second, compute
$(\mu_\aaa f\mu_\bb^*)(\mu_\cc g\mu_\dd^*)$ using $\mu_\bb^*\mu_\cc=\mu_{\cc'}\mu_{\bb'}^*$
with $\cc'=\cc/\gcd$, $\bb'=\bb/\gcd$, obtaining a term
$\mu_{\aaa\cc'}(\cdots)\mu_{\dd\bb'}^*$ whose two norms are equal but need not be $n$ or $m$.
\end{proof}

\begin{theorem}[The correct structure]\label{XIV:thm:structure}
$\cA^\sigma=\cA^{H}$ with $H=L^\perp\le\T^S$ the annihilator of the relation lattice, and
$\cA^\sigma$ carries a grading by $L$,
\[
\cA^\sigma=\overline{\bigoplus_{c\in L}}\ \cA^\sigma_{(c)},
\qquad
\cA^\sigma_{(c)}=\overline{\spn}\bigl\{\mu_\aaa f\mu_\bb^*:\aaa-\bb=c\bigr\},
\qquad
\cA^\sigma_{(0)}=C(Y_S),
\]
the grading being implemented by the residual gauge action of
$\T^S/H\cong\widehat L$. In particular $C(Y_S)=(\cA^\sigma)^{\widehat L}$.
\end{theorem}

\begin{proof}
$H=\overline{\iota(\R)}$ has annihilator $L$ by the proof of Theorem~\ref{XIV:thm:equiv}, so
$\T^S/H\cong\widehat L$ and the gauge action descends. The grading is the spectral
decomposition for that residual action, the degree of $\mu_\aaa f\mu_\bb^*$ being
$\aaa-\bb\in\Z^{(S)}$, which lies in $L$ precisely on $\cA^\sigma$. Degree zero is
Lemma~\ref{XIV:lem:gauge}.
\end{proof}

\begin{example}\label{XIV:ex:Qi}
Let $K=\Q(i)$ and $S=\{\pp,\bar\pp\}$ the two primes above $5$. Then $L=\Z(1,-1)$,
$\widehat L\cong\T$, and $\cA^\sigma$ is $\Z$-graded, the piece of degree $k(1,-1)$ being
spanned by the elements
\[
\mu_{\pp}^{m+k}\mu_{\bar\pp}^{\,n}\,f\,\bigl(\mu_{\pp}^{m}\mu_{\bar\pp}^{\,n+k}\bigr)^*,
\qquad\text{for instance}\qquad \mu_\pp^{k}f\mu_{\bar\pp}^{k\,*},
\]
both sides having the same norm $5^{m+n+k}$. The degree-zero part is $C(Y_S)$ and the
whole is a crossed-product-like algebra over it; it is non-commutative, and no spectral
argument using $\sigma$ alone can see inside it.
\end{example}

\section{Why the two proposed characterizations are circular}

\begin{remark}[Property A]\label{XIV:rem:A}
The proposal is to characterize $C(Y_S)$ as the unique MASA of $\cA^\sigma$ carrying a
faithful conditional expectation invariant under the gauge symmetries $\T^S$. By
Lemma~\ref{XIV:lem:gauge} this is correct as a \emph{description} --- $C(Y_S)$ is the range of the
averaging expectation $E_\gamma$ --- but it is not an intrinsic characterization inside
$(\cA,\sigma)$, because the $\T^S$-action is exactly the datum that $\sigma$ fails to
determine when $L\ne0$. By Theorem~\ref{XIV:thm:equiv}(3), $\sigma$ determines only
$H=L^\perp\subsetneq\T^S$. Invoking $\gamma$ to recover $C(Y_S)$ therefore assumes what is
missing.
\end{remark}

\begin{remark}[Property B]\label{XIV:rem:B}
The proposal is to characterize $C(Y_S)$ as the unique abelian MASA whose normalizer
generates $\cA_{K,S}$, i.e.\ as the unique Cartan subalgebra. That $C(Y_S)$ is \emph{a}
Cartan subalgebra is Theorem 2.5 of Chapter~\ref{ch:II} (topological principality; the
argument is reproduced inside the proof of Theorem~\ref{XIV:thm:uncond} below). That it is the
\emph{unique} one is precisely the statement needed to apply Renault's theorem \cite{Renault}
to a merely $\sigma$-equivariant isomorphism,
which is what Theorem 1.4 of Chapter~\ref{ch:XIII} was designed to supply and what is at issue here.
Assuming it is circular. We add that uniqueness of Cartan subalgebras is false for general
$C^*$-algebras and is a hard theorem where it holds, so it cannot be assumed lightly.
\end{remark}

\begin{remark}[What we do not know]\label{XIV:rem:open}
We have not found an intrinsic characterization of $C(Y_S)$ inside $(\cA_{K,S},\sigma)$ in the
unseparated case, and we have not found an obstruction either. The natural attempt --- to
exhibit a $\sigma$-commuting automorphism moving $C(Y_S)$, which would settle it negatively
--- fails for a reason worth recording. In a Cuntz algebra the quasi-free rotation
$\mu_i\mapsto\sum_ju_{ji}\mu_j$ is an automorphism because $\mu_i^*\mu_j=\delta_{ij}$; here
the isometries commute and satisfy $\mu_\pp^*\mu_\qq=\mu_\qq\mu_\pp^*$ for coprime
$\pp,\qq$, so, for $|\alpha|^2+|\beta|^2=1$ with $\alpha\beta\ne0$,
$(\alpha\mu_\pp+\beta\mu_\qq)^*(\alpha\mu_\pp+\beta\mu_\qq)=1+\bar\alpha\beta\,\mu_\qq\mu_\pp^*+\alpha\bar\beta\,\mu_\pp\mu_\qq^*\ne1$
and the rotated element is not an isometry. The Nica--Toeplitz relations are rigid enough to
block the obvious counterexample, which is a reason to think the question may have a positive
answer; we have not found it.
\end{remark}

\section{The gauge-equivariant repair}

\begin{definition}\label{XIV:def:gaugeequiv}
Call $\Phi:\cA_{K_1,S_1}\to\cA_{K_2,S_2}$ \emph{gauge-equivariant} if there is a topological
group isomorphism $\theta:\T^{S_1}\to\T^{S_2}$ with
$\Phi\circ\gamma^1_z=\gamma^2_{\theta(z)}\circ\Phi$ for all $z$. No bijection $S_1\to S_2$ is
presupposed, and none is encoded in $\theta$ by itself: the automorphisms of $\T^{S}$ are
dual to the automorphisms of the free abelian group $\Z^{(S)}$, most of which permute no
coordinates. Theorem~\ref{XIV:thm:uncond} produces a bijection, and shows that $\theta$ was
induced by one after all.
\end{definition}

\begin{theorem}[Unconditional reconstruction]\label{XIV:thm:uncond}
Let $K_1,K_2$ be arbitrary number fields and $S_1,S_2$ arbitrary infinite prime sets, and
let $\Phi:\cA_{K_1,S_1}\to\cA_{K_2,S_2}$ be a gauge-equivariant and $\sigma$-equivariant
isomorphism. Then $\Phi\bigl(C(Y_{S_1})\bigr)=C(Y_{S_2})$, and:
\begin{enumerate}
\item $\Phi$ is induced by an isomorphism $\theta_*$ of the underlying \'etale groupoids
which intertwines the scaling cocycles, $c_2\circ\theta_*=c_1$, and the integer-valued
gauge cocycles, $\widehat\theta\circ\widetilde c_2\circ\theta_*=\widetilde c_1$, where
$\widehat\theta:\Z^{(S_2)}\to\Z^{(S_1)}$ is the isomorphism dual to $\theta$;
\item $\theta_*$ induces a norm-preserving monoid isomorphism $\pi:J_{S_1}\to J_{S_2}$,
hence a norm-preserving bijection $S_1\to S_2$; moreover
$\widehat\theta(\delta_{\pi(\pp)})=\delta_\pp$, so $\theta$ is, a posteriori, the
isomorphism induced by the bijection $\pi$ of the coordinate circles;
\item the valuation fibres of $Y_{S_1}$ are carried homeomorphically onto those of
$Y_{S_2}$, compatibly with $\pi$; in particular the fibre over $0$, the compact space
$G_{S_1}$, is carried homeomorphically onto $G_{S_2}$;
\item the $\mathrm{KMS}_\beta$ simplices are affinely isomorphic for every $\beta$, so the
compact spaces $G_{S_i}/\Xi_\beta^\perp$ are homeomorphic, the orders $|\Xi_\beta|$ agree
when finite, the transition loci agree, and $\beta_c(S_1)=\beta_c(S_2)$;
\item if moreover $\Phi$ intertwines the symmetry actions along an isomorphism
$\iota:G_{S_1}\to G_{S_2}$, then the Frobenius classes are matched modulo inertia,
$\iota(\sigma_\pp)\in\sigma_{\pi(\pp)}\,r(\cO_{\pi(\pp)}^\times)$, and $\hat\iota^{-1}$
carries $\Xi_\beta(K_2,S_2)$ onto $\Xi_\beta(K_1,S_1)$ for every $\beta>0$.
\end{enumerate}
These are the conclusions of Theorem 2.5 of Chapter~\ref{ch:XIII}, with no norm-separation required
anywhere; what is conditional in (5) is conditional there too.
\end{theorem}

\begin{proof}
\emph{The Cartan pair.} By Lemma~\ref{XIV:lem:gauge}, $C(Y_{S_i})=\cA_i^{\T^{S_i}}$, and
gauge-equivariance, with $\theta$ surjective, gives
$\Phi(\cA_1^{\T^{S_1}})=\cA_2^{\T^{S_2}}$, i.e.\ $\Phi(C(Y_{S_1}))=C(Y_{S_2})$. Each pair
$\bigl(C(Y_S)\subset\cA_{K,S}\bigr)$ is Cartan with trivial Weyl twist: the groupoid
$\cG_S|_{Y_S}$ is \'etale, Hausdorff, second countable and amenable, and it is
topologically principal because the isotropy is trivial at every point with all valuations
finite ($\aaa\bb^{-1}y=y$ forces $\aaa=\bb$ at such a point) and those points are dense ---
this reproduces Theorem 2.5(1) of Chapter~\ref{ch:II}. Renault's theorem \cite{Renault}
therefore provides an isomorphism $\theta_*:\cG_{S_1}|_{Y_{S_1}}\to\cG_{S_2}|_{Y_{S_2}}$ of
\'etale groupoids inducing $\Phi$.

\emph{Both cocycles are intertwined.} For $f$ supported on an open bisection $B$, $\Phi(f)$
is supported on $\theta_*(B)$ (up to the $\T$-valued normalization of the trivialized
twist, which cancels below), and $\gamma^2_w$ multiplies it there by $w^{\widetilde
c_2}$. Comparing $\Phi(\gamma^1_z f)=\gamma^2_{\theta(z)}\Phi(f)$ pointwise gives
$z^{\widetilde c_1(\gamma)}=\theta(z)^{\widetilde c_2(\theta_*\gamma)}$ for all
$z\in\T^{S_1}$, i.e.\ $\widetilde c_1=\widehat\theta\circ\widetilde c_2\circ\theta_*$;
comparing $\Phi(\sigma^1_t f)=\sigma^2_t\Phi(f)$ gives $c_2\circ\theta_*=c_1$.

\emph{The semigroup, without norm-separation.} A full bisection of $\cG_S|_{Y_S}$ on which
$\widetilde c$ is constant contains, for each $y\in Y_S$, exactly one arrow $(g,y)$ with
$g=\aaa\bb^{-1}$ independent of $y$, $\aaa,\bb\in J_S$ coprime; that arrow exists only for
$y\in\bb Y_S$, and fullness forces $\bb Y_S=Y_S$, i.e.\ $\bb=1$. So these bisections are
exactly the $B_\aaa=\{(\aaa,y):y\in Y_S\}$, with $B_\aaa B_\bb=B_{\aaa\bb}$. This is
Lemma 2.1 of Chapter~\ref{ch:XIII} run with $\widetilde c$ in place of the scaling cocycle: there the
value of $c$ had to determine the element $g$, which is exactly norm-separation; the value
of $\widetilde c$ \emph{is} $g$, and no hypothesis is needed. Since $\widehat\theta$ is
injective, $\widetilde c_2\circ\theta_*$ is constant on a bisection iff $\widetilde c_1$
is, so $\theta_*$ maps the $B_\aaa$ to the $B_{\pi(\aaa)}$ for a monoid isomorphism
$\pi:J_{S_1}\to J_{S_2}$, which carries the free generators of the free abelian monoid
$J_{S_1}$ to those of $J_{S_2}$: a bijection $S_1\to S_2$. It preserves norms because
$\log\Ns\pi(\aaa)=c_2(B_{\pi(\aaa)})=c_1(B_\aaa)=\log\Ns\aaa$. Evaluating the intertwining
of $\widetilde c$ on $B_\pp$ gives
$\delta_\pp=\widetilde c_1(B_\pp)=\widehat\theta\bigl(\widetilde
c_2(B_{\pi(\pp)})\bigr)=\widehat\theta(\delta_{\pi(\pp)})$: $\widehat\theta$ maps standard
generators to standard generators, so $\theta$ is induced by $\pi$.

\emph{Fibres, simplices, symmetries.} $\theta_*$ maps $B_\aaa$ to $B_{\pi(\aaa)}$, hence
the range $\aaa Y_{S_1}$ to $\pi(\aaa)Y_{S_2}$, hence the valuation fibres to the
valuation fibres, and the fibre over $0$, which is $G_S$, homeomorphically onto its
counterpart. An isomorphism of $C^*$-dynamical systems carries $\mathrm{KMS}_\beta$ states
to $\mathrm{KMS}_\beta$ states; by [Ch.~\ref{ch:Main}, Thm.~3.6] the simplices are
$\operatorname{Prob}(G_S/\Xi_\beta^\perp)$, an affine isomorphism of Bauer simplices is a
homeomorphism of the extreme boundaries, and the transition locus is the set of $\beta$ at
which the simplex changes. Finally (5) is part (6) of Theorem 2.5 of Chapter~\ref{ch:XIII}: its proof
uses only the matched bisections and the torsor structure of the fibres, and nowhere uses
norm-separation, which entered that paper only through its Theorem 1.4 --- replaced here by
gauge-equivariance and Lemma~\ref{XIV:lem:gauge} --- and its Lemma 2.1 --- replaced by the
$\widetilde c$ argument above; its Remark 2.6 records that these are the only two entry
points.
\end{proof}

\begin{corollary}[Dedekind zeta]\label{XIV:cor:zeta}
If $S_i=P_{K_i}$ and $\Phi$ is gauge- and $\sigma$-equivariant, then the norm multisets
$\{\Ns\pp\}_{\pp\in P_{K_1}}$ and $\{\Ns\qq\}_{\qq\in P_{K_2}}$ agree with multiplicities,
hence
\[
\zeta_{K_1}(s)=\prod_{\pp}\bigl(1-\Ns\pp^{-s}\bigr)^{-1}=\zeta_{K_2}(s),
\]
so $K_1$ and $K_2$ are arithmetically equivalent. The narrow class numbers are \emph{not}
among the conclusions: $h^+_S$ is recovered by no argument of Chapter~\ref{ch:XIII} (its Remark 3.2),
and nothing here adds one.
\end{corollary}

\begin{remark}[Is the hypothesis reasonable?]\label{XIV:rem:reasonable}
Gauge-equivariance is a genuine strengthening and should be declared as such. It is however
natural: the gauge action is part of the standard data of a semigroup crossed product, it is
what makes $C(Y_S)$ the coefficient algebra, and any isomorphism arising from an isomorphism
of the underlying dynamical systems $(Y_{S_1},J_{S_1})\to(Y_{S_2},J_{S_2})$ is automatically
gauge-equivariant. What Theorem~\ref{XIV:thm:uncond} does not do is deduce the conclusion from
$\sigma$-equivariance alone, and by Corollary~\ref{XIV:cor:nothing} no argument through
$\cA^\sigma$ ever will.
\end{remark}

\section{Assessment}

\[
\begin{array}{lll}
\text{Statement} & \text{Status} & \text{Reference}\\\hline
C(Y_S)=\cA^{\T^S}\ \text{unconditionally} & \text{true} & \text{Lem.~1.1}\\
\text{norm-separated}\iff\cA^\sigma=C(Y_S) & \textbf{equivalent, not sufficient} & \text{Thm.~1.3}\\
\text{the hypothesis can be removed} & \textbf{false} & \text{Cor.~1.4}\\
P_K\ \text{norm-separated for}\ K\ne\Q & \textbf{never} & \text{Prop.~1.5}\\
\cA^\sigma=\bigoplus_n\cA^\sigma_n & \textbf{false: not a direct sum} & \text{Prop.~2.1}\\
\cA^\sigma\ \text{is}\ L\text{-graded},\ \cA^\sigma_{(0)}=C(Y_S) & \text{true} & \text{Thm.~2.2}\\
\text{Property A characterizes}\ C(Y_S) & \text{circular: uses}\ \gamma & \text{Rem.~3.1}\\
\text{Property B characterizes}\ C(Y_S) & \text{circular: assumes uniqueness} & \text{Rem.~3.2}\\
\text{intrinsic recovery from}\ (\cA,\sigma) & \textbf{open} & \text{Rem.~3.3}\\
\text{gauge-equivariant reconstruction} & \textbf{true, all}\ K,S & \text{Thm.~4.2}\\
\theta\ \text{induced by a bijection}\ S_1\to S_2 & \text{true a posteriori, not a priori} & \text{Thm.~4.2(2)}\\
\text{Frobenius classes, }\{\Xi_\beta\}\ \text{matched} & \text{if symmetries intertwined} & \text{Thm.~4.2(5)}\\
\zeta_{K_1}=\zeta_{K_2}\ (S_i=P_{K_i}) & \text{true, gauge-equivariantly} & \text{Cor.~4.3}\\
h^+_{K_1}=h^+_{K_2} & \textbf{not claimed} & \text{Cor.~4.3}\\
\end{array}
\]

The question was well posed and its answer is a change of shape. Norm-separation cannot be
removed from Theorem 1.4 of Chapter~\ref{ch:XIII}, because Theorem~\ref{XIV:thm:equiv} shows it is that
theorem's conclusion in disguise; what can be done is to replace $\sigma$-equivariance by
gauge-equivariance, whereupon the reconstruction holds with the exact strength of Paper
XIII's Theorem 2.5 for every number field and every prime set, including the classical case
$S=P_K$ that Proposition~\ref{XIV:prop:PK} shows was previously excluded --- with what was
conditional there (the Frobenius classes, when the symmetry actions are not known to be
intertwined) still conditional, and the narrow class number still not recovered.

Two questions remain, and the first is the sharp form of the original one. Is $C(Y_S)$
intrinsically characterized inside $(\cA_{K,S},\sigma)$ when $L\ne0$? Remark~\ref{XIV:rem:open}
shows the obvious counterexample is blocked by the Nica--Toeplitz relations, so we suspect
the answer is yes and that the mechanism is the rigidity of commuting isometries rather than
any spectral argument. Second: the $L$-graded algebra $\cA^\sigma$ of
Theorem~\ref{XIV:thm:structure} is a new object in this programme, being the first non-commutative
fixed-point algebra to appear; whether its own structure --- for instance whether it is
simple, or its ideal lattice --- encodes the relation lattice $L$, and hence the pattern of
coincidences among the norms, we have not examined.

